\subsection{解析函数}

3.2.1. 设 U 是 E 的开子集，f 是从 U 到 F 的映射。若对 U 中每一点 a，都存在收敛级数 \(C^{\omega}\)，使得当 E 中的 x 充分接近零时 \(f_{a} \in \mathcal{H}(E; F)\)，则称 f 在 U 中属于 \(f(a + x) = f_{a}(x)\) 类，或称 K-解析（或简写为解析）。当 K = R（分别地，C）时，也称 f 为实解析（分别地，复解析或全纯）。从 U 到 F 的解析映射构成所有从 U 到 F 的映射空间的向量子空间，记作 \(\mathcal{C}^{\omega}(U; F)\)。

对于 \(a \in \mathbf{U}\)，形式级数 \(f_{a}\) 是唯一的：称其为 \(f\) 在点 \(a\) 处的幂级数展开。若 \(f_{a} = \sum_{\alpha}(f_{a})_{\alpha}\)（其中 \((f_{a})_{\alpha} \in \mathrm{P}_{\alpha}(\mathrm{E}_{1}, \ldots, \mathrm{E}_{n}; \mathrm{F})\)），置：

\[
\Delta^ {\alpha} f (a) = \left(f _ {a}\right) _ {\alpha}.
\]

3.2.2. 若 \(f \in \mathcal{C}^{\omega}(\mathrm{U};\mathrm{F})\)，则映射 \(\Delta^{\alpha}f: a \mapsto \Delta^{\alpha}f(a)\) 从 U 到

\[
\mathbf {P} _ {\alpha} (\mathrm{E} _ {1}, \dots , \mathrm{E} _ {n}; \mathrm{F})
\]

是解析的。映射 \(\Delta^{\alpha}: f \mapsto \Delta^{\alpha}f\) 是从 \(\mathcal{C}^{\omega}(\mathrm{U};\mathrm{F})\) 到 \(\mathcal{C}^{\omega}(\mathrm{U};\mathrm{P}_{\alpha}(\mathrm{E}_{1},\ldots,\mathrm{E}_{n};\mathrm{F}))\) 的 K-线性映射。因此，对于 \(a \in \mathrm{U}\) 以及 \(\alpha, \beta \in \mathbf{N}^{n}\)，有 \(\Delta^{\beta}(\Delta^{\alpha}f)(a) \in \mathrm{P}_{\beta}(\mathrm{E}_{1},\ldots,\mathrm{E}_{n};\mathrm{P}_{\alpha}(\mathrm{E}_{1},\ldots,\mathrm{E}_{n};\mathrm{F}))\)。若 \(x = (x_{i}) \in \mathrm{E}\)，则 \((\Delta^{\beta}(\Delta^{\alpha}f)(a))(x) \in \mathrm{P}_{\alpha}(\mathrm{E}_{1},\ldots,\mathrm{E}_{n};\mathrm{F})\)，且 \(((\Delta^{\beta}(\Delta^{\alpha}f)(a))(x))(x) \in \mathrm{F}\)。F 中此元素等于 \(((\alpha, \beta)) (\Delta^{\alpha + \beta}f(a))(x)\)。这表示为：

\[
\Delta^ {\beta} \circ \Delta^ {\alpha} = ((\alpha , \beta)) \Delta^ {\alpha + \beta}.
\]

\(n=1\)3.2.3. \(f\) 和 \(\Delta^{\alpha}f\) 在 U 中同一点 \(a\) 处的幂级数展开的严格收敛指标（以及当 n=1 时的严格收敛半径）相同。

3.2.4. 设 \(f \in \mathcal{C}^{\omega}(\mathrm{U};\mathrm{F})\)。则 \(f\) 在 \(\mathbf{U}\) 中严格可微且无穷次可微（若 \(\mathbf{K} = \mathbf{R}\)，则 \(f\) 在 \(\mathbf{C}^{\infty}\) 中属于 \(\mathbf{U} \) 类）。\(f\) 的迭代导数是解析的，且它们在点 \(a\) 处的值是对称多线性映射。于是可如第 2.4.2 号引入迭代偏导数的记号 \(\mathbf{D}^{a}f\)。有：

\[
\alpha ! \Delta^ {\alpha} f (a) (h) = D ^ {\alpha} f (a). (h, \dots , h)
\]

对所有 \(a \in \mathbf{U}\) 和 \(h \in \mathbf{E}\) 均成立，记作：

\[
\alpha ! \Delta^ {\alpha} f = D ^ {\alpha} f.
\]

3.2.5. 设 U 是 E 的开子集，f、g 是从 U 到 F 的两个解析映射。设 \(a \in U\)。f 和 g 在 a 处接触阶 \(\geqslant k\) 的充要条件是，对所有满足 \(\Delta^{\alpha}f(a) = \Delta^{\alpha}g(a)\) 的 \(\alpha\)，都有 \(|\alpha| \leqslant k\)。若 f 和 g 在 a 处具有无穷阶接触，则它们在 a 的某邻域上相等。U 中 f 和 g 具有无穷阶接触的点集是开且闭的。

特别地，若 U 连通，且存在 U 的一个非空开子集，在其上 f 和 g 相等，则 f = g（“解析延拓原理”）。

3.2.6. 设 U 是 E 的开子集，f 是从 U 到 F 的解析映射。若 f 的导数 Df 为零，则 f 局部为常值。

3.2.7. 假设 F 拟完备，且 G 是完备赋范空间。设 g 是从 E 的开子集 U 到 G 的解析映射，f 是从 G 的包含 g(U) 的开子集 V 到 F 的解析映射。复合映射 \(f \circ g\) 在 U 中解析。再设 \(0 \in U\) 且 \(g(0) = 0\)。将 f 在 0 处的幂级数展开中的变量替换为 g 在 0 处的幂级数展开，即得到 \(f \circ g\) 在 0 处的幂级数展开（3.1.9）。

3.2.8. 设 \(F_{1},\ldots,F_{m}\) 是分离多范数空间，u 是从 \(F_{1}\times\cdots\times F_{m}\) 到 F 的连续多线性映射。设 U 是 E 的开子集，且 \(f_{i}\in\mathcal{C}^{\omega}(U;F_{i})\)。函数 \(u(f_{1},\ldots,f_{m})\) 是解析的，且其在点 \(a\in U\) 处的幂级数展开是级数 \(u((f_{1})_{a},\ldots,(f_{m})_{a})\)（3.1.8）。

3.2.9. 假设 \(F\) 拟完备。设 \(f \in \mathcal{H}(E_1, \ldots, E_n; F)\)；函数 \(x \mapsto f(x)\)（3.1.7）在开集 \(C(f)\) 中解析，其中 C(\(f\)) 是 f 的严格收敛域。若 \(n = 1\) 且 \(\|a\| < \rho(f)\)，则 \(f\) 在 \(a\) 处的幂级数展开的严格收敛半径至少等于 \(\rho(f) - \|a\|\)。若 \(\rho(f) = +\infty\)，则称 \(f\) 为整函数。

3.2.10. 保持 3.2.9 的假设。若 K = C，则将 C(f) 换为 \(\tilde{C}(f)\)，并将 \(\rho(f)\)（当 \(\tilde{\rho}(f)\) 时）换为 \(n = 1\) 后，3.2.9 的结论仍完全成立。若 K = R，则函数 \(x \mapsto f(x)\) 在 \(\tilde{C}(f)\) 中解析。

3.2.11. 假设 \(E_{i} = K\)（\(1 \leqslant i \leqslant n\)）。设 U 是 \(K^{n}\) 的开子集，且 \(f \in \mathcal{C}^{\omega}(K^{n}; F)\)。若 \(0 \in U\)，且 \(f_{0} = \sum_{\alpha} X^{\alpha} c_{\alpha}\) 是 f 在 0 处的幂级数展开，则在将 \(\Delta^{\alpha} f\) 识别为 F 后，\(P_{\alpha}(K^{n}; F)\) 在 0 处的幂级数展开写作：

\[
(\Delta^ {\alpha} f) _ {0} = \sum_ {\beta} ((\alpha , \beta)) X ^ {\beta} c _ {\alpha + \beta}.
\]

特别地，对 \(1 \leqslant i \leqslant n\) 以及 \(x \in \mathbf{C}(f_0)\)：

\[
\partial_ {i} f (x) = \sum_ {\alpha} (\alpha_ {i} + 1) x ^ {\alpha} c _ {\alpha + \varepsilon_ {i}}.
\]

